\documentclass[a4paper,12pt,fleqn]{article}%fleqn pour aligner les equations à gauche
\usepackage{../../Styles/Eval}


\newcommand{\chnum}{Chapitre 5}
\newcommand{\chtitle}{Cortège électronique et stabilité}
\newcommand{\dtype}{DS 3}

\begin{document}
\maketitle
\section{Cours}
\begin{enumerate}
    \item Donner une définition de la configuration électronique.
\end{enumerate}
\noindent\fbox{
    \begin{minipage}{\textwidth}
        \vspace{5em}
        \hspace{\linewidth}
    \end{minipage}
}


\begin{enumerate}[resume]
    \item Expliquer la règle de formation des ions monoatomiques.
\end{enumerate}

\noindent\fbox{
    \begin{minipage}{\textwidth}
        \vspace{8em}
        \hspace{\linewidth}
    \end{minipage}
}
\section{Configuration électronique et tableau périodique}
\begin{center}
    \begin{tabular}{|m{.1\linewidth}|m{.2\linewidth}|m{.15\linewidth}|m{.1\linewidth}|m{.1\linewidth}|m{.15\linewidth}|}
        \hline
        Atome & Configuration électronique & Nombre d'électrons de valence & Ligne & Colonne & Ion monoatomique formé\\
        \hline
        \ce{Si} (Z = 14) & & & & & \\
        \hline
        \ce{Li} (Z = 3) & & & & & \\
        \hline
        \ce{Cl} (Z = 17) & & & & & \\
        \hline
        \ce{Be} (Z = 4) & & & & & \\
        \hline            
    \end{tabular}
\end{center}
\begin{enumerate}[resume]
    \item Compléter les colonnes ["Configuration électronique"] et ["Nombre d'électrons de valence"] du tableau.
    \item Expliquer comment sont placés les éléments dans la tableau périodique, puis compléter les colonnes ["Ligne"] et ["Colonne"] du tableau.
\end{enumerate}
\noindent\fbox{
    \begin{minipage}{\textwidth}
        \vspace{5em}
        \hspace{\linewidth}
    \end{minipage}
}
\begin{enumerate}[resume]
    \item Compléter la colonne ["Ion monoatomique formé"]
\end{enumerate}
\section{Stabilité des molécules}
\begin{enumerate}
    \item Expliquer la règle de stabilité associée aux molécules.
\end{enumerate}
\noindent\fbox{
    \begin{minipage}{\textwidth}
        \vspace{8em}
        \hspace{\linewidth}
    \end{minipage}
}
\begin{enumerate}[resume]
    \item Analyser la stabilité des molécules suivantes (annotations sur le schéma en vert). Corriger les erreurs éventuelles (annotation sur le schéma en rouge). Justifier les corrections.
\end{enumerate}
\begin{center}
    \begin{tabular}{|>{\centering}m{.48\linewidth}|>{\centering\arraybackslash}m{.48\linewidth}|}
        \hline
        \chemfig{H-C(-[2]H)(-[6]H)-O-H} & \chemfig{H-N=O}\\
        % \hline
        \vspace{4cm}. & \vspace{4cm}.\\
        \hline
        \chemfig{N-N} & \chemfig{H-O-C-N(-[2]H)-H}\\
        % \hline
        \vspace{4cm}. & \vspace{4cm}.\\
        \hline
    \end{tabular}
\end{center}
\end{document}